\noindent
Audit findings are both outcomes of prior work and inputs to future audit
planning. When prior findings change subsequent verification, the finding
record combines information about underlying conditions with information about
where auditors looked. We characterize what can be learned from that record.
Three programs whose underlying condition is persisting, unchanged, or improving
can produce the same finding history, so the record alone does not identify the
underlying association. It does bound it: if follow-up can only raise detection,
the observed finding risk ratio is an upper bound on the underlying-condition
risk ratio, and a defensible ceiling on the follow-up detection multiplier makes
the bound two-sided. Reference observation under a common measurement rule, or
other information about detection, can make the bounds operational. Prior
findings are natural routing information for follow-up; what they reveal about
discovery cannot be inferred from the finding record alone. In Federal Audit
Clearinghouse data, prior findings predict both recurrence and materially more
newly recorded findings among programs audited in depth in adjacent years,
while a recurrence-share statistic ranges from 61.2 to 90.9 percent on a
reporting convention alone. A risk score validated against recorded findings
can improve while carrying no information about underlying conditions.
